\subsection{The order function}

Let A be a filtered ring, E a filtered A-module and \((\mathrm{E}_{n})\) the filtration of E. For all \(x \in E\) let \(v(x)\) denote the least upper bound in \(\bar{R}\) of the set on integers \(n \in Z\) such that \(x \in E_{n}\) . Then the following equivalences hold:

\[
\left\{ \begin{array}{l} v (x) = - \infty \Leftrightarrow x \notin \bigcup_ {n \in \mathbf {Z}} \mathrm{E} _ {n} \\ v (x) = p \quad \Leftrightarrow x \in \mathrm{E} _ {p} \quad \text { and } \quad x \notin \mathrm{E} _ {p + 1} \\ v (x) = + \infty \Leftrightarrow x \in \bigcap_ {n \in \mathbf {Z}} \mathrm{E} _ {n} \end{array} \right.\tag{4}
\]

The mapping \(v: \mathrm{E} \to \overline{\mathbf{R}}\) is called the order function of the filtered module \(\mathrm{E}\) . If \(v\) is known then so are the \(\mathrm{E}_n\) , for \(\mathrm{E}_n\) is the set of \(x \in \mathrm{E}\) such that \(v(x) \geqslant n\) ; the fact that the \(\mathrm{E}_n\) are additive subgroups of \(\mathrm{E}\) implies the relation

\[
v (x - y) \geqslant \inf (v (x), v (y)).\tag{5}
\]

The above definition applies in particular to the filtered A-module \(A_{s}\) ; let w be its order function. It follows from equation (3) of no. 1 that for \(a \in A\) and \(x \in E\) ,

\[
v (a x) \geqslant w (a) + v (x)\tag{6}
\]

whenever the right-hand side is defined; in particular, for \(a \in A\) and \(b \in A\) ,

\[
w (a b) \geqslant w (a) + w (b)\tag{7}
\]

whenever the right hand side is defined.

The order function is defined similarly on a filtered group G which is not necessarily commutative; the corresponding relation to (5) is then written

\[
v (y x ^ {- 1}) = v (x y ^ {- 1}) \geqslant \inf (v (x), v (y)).\tag{5'}
\]

